\documentclass[12pt, letterpaper]{article}
\usepackage{graphicx}
\usepackage[margin=1in]{geometry}
\usepackage{amsmath, amssymb, amsthm, mathrsfs}
\usepackage{fancyhdr}
\usepackage{tikz}
\usetikzlibrary{shapes.geometric}
\usepackage{enumerate}
\usepackage{cancel}
\usetikzlibrary{positioning}
\usetikzlibrary{calc}
\usepackage[utf8]{inputenc}
\usepackage{xcolor}
\usepackage{array}
\usepackage{empheq}
\usepackage{framed}

\title{HW 1.1}
\author{Your Name}
\date{\today}
\pagestyle{fancy}
\setlength{\headheight}{42pt}
\renewcommand{\headrulewidth}{0pt}
\renewcommand{\footrulewidth}{0pt}
\fancyhf{}
\rhead{
    Zack Abdirashid\\
    4/17/26 \\
    MATH 402
}
\rfoot{Page \thepage}

\begin{document}
\paragraph{\underline{10} \\ \\}

\textbf{1.} \begin{tabular}[t]{@{}p{0.93\textwidth}@{}}
    (a) What is the smallest positive integer $n$ such that there are exactly four non-isomorphic Abelian groups of order $n$? Name the four groups. \\[6pt] \\
    \begin{tabular}[t]{l}
        By the Fundamental Theorem of Finite Abelian Groups, we want to find a prime $p$ such \\
        that the number of partitions of $p$ is 4. Consider the product $p^{2}_{1}p^{2}_{2}$ where $p_{1}$ and $p_{2}$ are di-\\
        stinct primes. Since both amount to 2 partitions respectively, their product amounts to 4 \\
        partitions. To get the smallest integer $n$, we want to use the smallest primes that create \\
        4 partitions. So,
    \end{tabular} 
    \begin{align*}
        n = 2^{2} \cdot 3^{2} = \boxed{36}
    \end{align*}
    \begin{tabular}[t]{l}
        The groups should be of the forms $\mathbb{Z}_{p_1^{2}} \oplus \mathbb{Z}_{p_2^{2}}, \mathbb{Z}_{p_1^{2}} \oplus (\mathbb{Z}_{p_2} \oplus \mathbb{Z}_{p_2}), (\mathbb{Z}_{p_1} \oplus \mathbb{Z}_{p_{1}}) \oplus \mathbb{Z}_{p_2^{2}}, $ \\
        and $(\mathbb{Z}_{p_1} \oplus \mathbb{Z}_{p_1}) \oplus (\mathbb{Z}_{p_2} \oplus \mathbb{Z}_{p_2})$. Thus, the groups are
        \end{tabular} 
        \begin{center}
        \fbox{\begin{tabular}{r@{\quad}l}
            1. & $\mathbb{Z}_{4} \oplus \mathbb{Z}_{9}$ \\
            2. & $\mathbb{Z}_{4} \oplus (\mathbb{Z}_{3} \oplus \mathbb{Z}_{3})$ \\
            3. & $(\mathbb{Z}_{2} \oplus \mathbb{Z}_{2}) \oplus \mathbb{Z}_{9}$ \\
            4. & $(\mathbb{Z}_{2} \oplus \mathbb{Z}_{2}) \oplus (\mathbb{Z}_{3} \oplus \mathbb{Z}_{3})$
        \end{tabular}}
        \end{center}
    (b) Up to isomorphism, how many additive Abelian groups of order 16 have the property that $x + x + x + x = 0$ for all $x$ in the group? \\[6pt]
    \begin{tabular}[t]{l}
        Since 2 is prime and $2^{4} = 16$, we want to find the groups corresponding to the partitions of $p^{4}$. \\
        Since partitions of $p^{4}$ are $p^{4}, p^{3}p, p^{2}p^{2}, p^{2}pp$, and $pppp$, our choices of groups are \\
        \end{tabular}
        \begin{center}
        {\begin{tabular}{r@{\quad}l}
            &1. $\mathbb{Z}_{16}$ \ \text{(max order of 16)} \\
           &2. $\mathbb{Z}_{8} \oplus \mathbb{Z}_{2}$ \ \text{(max order of 8)} \\
            &3. $\mathbb{Z}_{4} \oplus \mathbb{Z}_{4}$ \ \text{(max order of 4)} \\
            &4. $\mathbb{Z}_{4} \oplus (\mathbb{Z}_{2} \oplus \mathbb{Z}_{2})$ \ \text{(max order of 4)} \\
            &5. $\mathbb{Z}_{2} \oplus \mathbb{Z}_{2} \oplus \mathbb{Z}_{2} \oplus \mathbb{Z}_{2}$ \ \text{(max order of 2)}
            \end{tabular}}
        \end{center}
        \begin{tabular}[t]{l}
            Groups whose elements have a max order of 4 satisfy the property. Also, groups with max order\\
            of 2 also satisfy it since $(x + x) + (x + x) = 0 + 0 = 0$. So, \boxed{\text{3 groups satisfy the property.}}
            \end{tabular} \\ \\
    (c) Let $p_1, p_2, \dots, p_n$ be distinct primes. Up to isomorphism, how many Abelian groups are there of order $p_1^4 p_2^4 \cdots p_n^4$? \\[6pt]
    \begin{tabular}[t]{l}
        For each $p_{i}^4$, it has 5 partitions. So, the total number of groups of this order is \boxed{5^{n}}
        \end{tabular}
\end{tabular}

\bigskip

\textbf{2.} \begin{tabular}[t]{@{}p{0.93\textwidth}@{}}
    Let $G = \{1, 7, 43, 49, 51, 57, 93, 99, 101, 107, 143, 149, 151, 157, 193, 199\}$ under multiplication modulo 200. Determine the isomorphism class of $G$. \textit{Justify your answer.}
\end{tabular}
\indent \begin{tabular}[t]{l}
    Since $|G| = 16$, we want to find E.D.Ps of order 16. Since $2^4 = 16$, the only options are 
    \end{tabular}
    \begin{alignat*}{2}
        & \mathbb{Z}_{16} &\qquad& \text{possible orders: } 1, 2, 4, 8, 16 \\
        & \mathbb{Z}_{8} \oplus \mathbb{Z}_{2} && \text{possible orders: } 1, 2, 4, 8 \\
        & \mathbb{Z}_4 \oplus \mathbb{Z}_4 && \text{possible orders: } 1, 2, 4 \\
        & \mathbb{Z}_4 \oplus (\mathbb{Z}_2 \oplus \mathbb{Z}_2) && \text{possible orders: } 1, 2, 4 \\
        & \mathbb{Z}_2 \oplus \mathbb{Z}_2 \oplus \mathbb{Z}_2 \oplus \mathbb{Z}_2 && \text{possible orders: } 1, 2
    \end{alignat*}
    \indent \begin{tabular}[t]{l}
    The orders of the elements in $G$ under multiplication modulo 200 are:
    \end{tabular}
    \begin{gather*}
        |1| = 1, \quad |7| = 4, \quad |43| = 4, \quad |49| = 2, \\
        |51| = 2, \quad |57| = 4, \quad |93| = 4, \quad |99| = 2, \\
        |101| = 2, \quad |107| = 4, \quad |143| = 4, \quad |149| = 2, \\
        |151| = 2, \quad |157| = 4, \quad |193| = 4, \quad |199| = 2
    \end{gather*}
    \indent \begin{tabular}[t]{l}
        Since no elements have order 8 or 16, we can eliminate $\mathbb{Z}_{16}$ and $\mathbb{Z}_{8} \oplus \mathbb{Z}_{2}$. Also, since there's at \\
        least one element of order 3, we can eliminate $\mathbb{Z}_2 \oplus \mathbb{Z}_2 \oplus \mathbb{Z}_2 \oplus \mathbb{Z}_2$. So, the only options left \\
        are \boxed{\text{$\mathbb{Z}_4 \oplus \mathbb{Z}_4$ and $\mathbb{Z}_4 \oplus (\mathbb{Z}_2 \oplus \mathbb{Z}_2)$}}\\
    \end{tabular}

\bigskip

\textbf{3.} \begin{tabular}[t]{l}
    Prove Theorem 12.1, which says the following:
\end{tabular}

\begin{quote}
    Let $a$, $b$, and $c$ belong to a ring $R$. Then
    \begin{enumerate}
        \item $a0 = 0a = 0$
        \begin{proof}
            $a0 = a(0 + 0) = a0 + a0$. Adding $(-a0)$ to both sides,
            \begin{align*}
                a0 + (-a0) &= (a0 + a0) + (-a0) \\
                \Rightarrow 0 &= a0 
                \end{align*}
            Similarly, 
            \begin{align*}
                0a &= (0 + 0)a = 0a + 0a \\
                \Rightarrow 0a + (-0a) &= (0a + 0a) + (-0a) \\
                \Rightarrow 0 &= 0a 
                \end{align*}
        \end{proof}
        \item $a(-b) = (-a)b = -(ab)$
        \begin{proof}
            If a(-b) = -(ab), then that implies a(-b) is the additive inverse of ab. So,
            \begin{align*}
                a(-b) + ab &= a(-b + b) = a0 = 0.   
                \end{align*}
                So,
                \begin{align*}
                    a0 + (-(ab)) &= 0 + (-(ab)) \\
                    \Rightarrow a(0 + (-b)) &= -(ab) \\
                    \Rightarrow a(-b) &= -(ab)
                    \end{align*}
            Similarly,
            \begin{align*}
                (-a)b + ab &= (-a + a)b = 0b = 0.  \\
                \Rightarrow 0b + (-(ab)) &= 0 + (-(ab)) \\
                \Rightarrow (0 + (-a))b &= -(ab) \\
                \Rightarrow (-a)b &= -(ab)
                \end{align*}
        \end{proof}
        \item $(-a)(-b) = ab$
        \begin{proof}
            By Property 2 of Theorem 12.1,
            \begin{align*}
                (-a)(-b) &= -((-a)b) \\
                &= -(-(ab)) \\
                &= ab
                \end{align*}
        \end{proof}
        \item $a(b - c) = ab - ac$ and $(b - c)a = ba - ca$
        \begin{proof}
            By the definition of rings and Property 2 of Theorem 12.1,
            \begin{align*}
                a(b - c) &= a(b + (-c)) \\
                &= ab + a(-c) \\
                &= ab + (-(ac)) \\
                &= ab - ac
            \end{align*}
            Similarly,
            \begin{align*}
                (b - c)a &= (b + (-c))a \\
                &= ba + (-ca) \\
                &= ba - ca
                \end{align*}
        \end{proof}
    \end{enumerate}
    Furthermore, if $R$ has a unity element 1, then
    \begin{enumerate}
        \setcounter{enumi}{4}
        \item $(-1)a = -a$
        \begin{proof}
            By Property 2 of Theorem 12.1,
            \begin{align*}
                (-1)a &= -(1a) \\
                &= -(a) \\
                &= -a
                \end{align*}
        \end{proof}
        \item $(-1)(-1) = 1$
        \begin{proof}
            By Property 3 of Theorem 12.1,
            \begin{align*}
                (-1)(-1) &= (1)(1) \\
                &= 1
                \end{align*}
        \end{proof}
    \end{enumerate}
\end{quote}

\bigskip

\textbf{4.} \begin{tabular}[t]{@{}p{0.93\textwidth}@{}}
    (a) Find an integer $n$ so that the ring $\mathbb{Z}_n$ does \textbf{not} satisfy any of the following properties for all $a, b, c \in \mathbb{Z}_n$:
\end{tabular}

\begin{quote}
    \begin{enumerate}[i.]
        \item If $a^2 = a$, then $a = 0$ or $a = 1$ 
        \item If $ab = 0$, then $a = 0$ or $b = 0$
        \item If $ab = ac$ and $a \neq 0$, then $b = c$.
    \end{enumerate}
\end{quote} 

\indent \indent \begin{tabular}[t]{l} 
    Let $n = 6$ and consider the ring $\mathbb{Z}_6$:  \\
    \indent i. $3^{2} = 9 \ \text{mod} \ 6 = 3$, but $3 \neq 0$ and $3 \neq 1$. \\
    \indent ii. $(3)(2) = 6 \ \text{mod} \ 6 = 0$, but $3 \neq 0$ and $2 \neq 0$. \\
    \indent iii. $(3)(4) = 12 \ \text{mod} \ 6 = 0$ and $(3)(2) = 6 \ \text{mod} \ 6 = 0$, but $4 \neq 2$. \\
    \end{tabular} \\
\begin{tabular}[t]{@{}p{0.93\textwidth}@{}}
    \phantom{\textbf{4.} }(b) Show that the three properties in part (a) are satisfied in $\mathbb{Z}_p$, where $p$ is prime.
\end{tabular} \\

\begin{tabular}[t]{l}
    Since $\mathbb{Z}_p$ has a unity of 1, it has multiplicative inverses. So, \\
    \indent i.
    \end{tabular}
    \begin{align*}
        a^{2}(a^{-1}) &= a(a^{-1}) \\
        \Rightarrow a &= 0.
        \end{align*}
        \indent \indent \begin{tabular}[t]{l}
            If $a \neq 0$, then
            \end{tabular}
            \begin{align*}
                a^{2} &= a \\
                \Rightarrow a^2 - a &= 0 \\
                \Rightarrow a(a - 1) &= 0 \\
               \Rightarrow a^{-1}(a(a - 1)) &= a^{-1}0 \\
               \Rightarrow a - 1 & = 0 \\
               \Rightarrow a &= 1
                \end{align*}
    \indent \indent \begin{tabular}[t]{l}
        ii.
        \end{tabular}
        \begin{align*}
            ab(b^{-1}) &= 0(b^{-1}) \\
            \Rightarrow a &= 0 \ \\
            &\text{or}  \\ 
            a^{-1}(ab) &= a^{-1}0 \\
            \Rightarrow b &= 0
            \end{align*}
   \indent \indent \begin{tabular}[t]{l}
    iii.
        \end{tabular}
        \begin{align*}
            (a^{-1})(ab) &= (a^{-1})(ac) \\
            \Rightarrow b &= c
            \end{align*}

\bigskip

\textbf{5.} \begin{tabular}[t]{@{}p{0.93\textwidth}@{}}
    (a) Give an example of ring elements $a$ and $b$ with the property that $ab = 0$, but $ba \neq 0$. \\[6pt]
    (b) Let $n$ be an integer greater than 1. In a ring in which $x^n = x$ for all $x$, show that $ab = 0$ implies that $ba = 0$.
\end{tabular} \\

\indent \indent \begin{tabular}[t]{l}
    (a) If $ab = ba = 0$, that means $a$ and $b$ commute. So, we want to find ring elements that \\
    don't commute. Consider the ring $M_2(\mathbb{Z})$ and let $A = \begin{bmatrix} 0 & 1 \\ 0 & 0 \end{bmatrix}$ and $B = \begin{bmatrix} 1 & 0 \\ 0 & 0 \end{bmatrix}$ be elements \\
    of the ring. Then,
    \end{tabular}
    \begin{align*}
        AB &= \begin{bmatrix} 0 & 1 \\ 0 & 0 \end{bmatrix} \begin{bmatrix} 1 & 0 \\ 0 & 0 \end{bmatrix} = \begin{bmatrix} 0 & 0 \\ 0 & 0 \end{bmatrix} \\
        BA &= \begin{bmatrix} 1 & 0 \\ 0 & 0 \end{bmatrix} \begin{bmatrix} 0 & 1 \\ 0 & 0 \end{bmatrix} = \begin{bmatrix} 0 & 1 \\ 0 & 0 \end{bmatrix}
    \end{align*}
    \indent \indent \indent So, $AB \neq BA$. \\\\
\indent \indent \begin{tabular}[t]{l}
    (b) Let $n = 2$. Then, 
    \end{tabular}
    \begin{align*}
        (ba)^{2} &= ba(ba) \\
        &= b(ab)a \\
        &= b(0)a \\
        &= 0
        \end{align*}
    \indent \indent \indent By the definition of the ring, $(ba)^{2} = ba$. So, $ba = 0$.
\bigskip

\textbf{6.} \begin{tabular}[t]{@{}p{0.93\textwidth}@{}}
    Let $a$ be a fixed element of a ring $R$, and let $S = \{x \in R \mid ax = 0\}$. Show that $S$ is a subring of $R$.
\end{tabular}
\begin{proof}
    Let $R$ be a ring, $a$ be a fixed element of $R$, and $S = \{x \in R \mid ax = 0\}$. By the Subring Test, we need to show that $S$ is closed under subtraction and multiplication. Take $x, y \in S$ and consider the product $a(xy)$. By Theorem 12.1,
    \begin{align*}
        a(xy) &= (ax)y \\
        &= 0y \\
        &= 0
        \end{align*}
    So, $xy \in S$. Similarly,
    \begin{align*}
        a(x-y) &= ax - ay \\
        &= 0 - 0 \\
        &= 0.
    \end{align*}
    So, $x-y \in S$. By the Subring Test, $S$ is a subring of $R$.
    \end{proof}

\end{document}